\documentclass[12pt]{article}
\usepackage[a4paper,margin=1in]{geometry}
\usepackage[utf8]{inputenc}
\usepackage[T1]{fontenc}
\usepackage{amsmath,amssymb}
\usepackage{textcomp}
\usepackage{parskip}
\begin{document}
\begin{center}
{\Large \textbf{2012 STEP III}}
\end{center}

\subsection*{Page 1}
\begin{flushleft}
\bigskip\textbf{Section A: Pure Mathematics}\par
dy\par
n\par
\bigskip\textbf{Question 1.} Given that z = y , show that\par
dx\par
\lambda \mu\par
\[dz dy dy 2 d2y\]
n-1\par
\[= y n + 2y .\]
\[dx dx dx dx2\]
!\par
\lambda \mu\par
\textbf{(i)} Use the above result to show that the solution to the equation\par
\[dy 2 d2y\]
\[+ 2y = \sqrt\{y\} (y > 0)\]
\[dx dx2\]
\lambda \mu\par
dy 2\par
\[that satisfies y = 1 and = 0 when x = 0 is y = 3 x^2 + 1 3 .\]
dx 8\par
\textbf{(ii)} Find the solution to the equation\par
\[dy 2 d2y\]
\[y + y = 0\]
\[dx - dx2\]
\lambda \mu\par
dy\par
\[that satisfies y = 1 and = 0 when x = 0.\]
dx\par
\bigskip\textbf{Question 2.} In this question, x < 1 and you may ignore issues of convergence.\par
| |\par
\textbf{(i)} Simplify\par
2 4 2n\par
(1 x)(1 + x)(1 + x )(1 + x ) (1 + x ) ,\par
- \cdots\par
where n is a positive integer, and deduce that\par
1 2 4 2n\par
\[x2n+1\]
\[= (1 + x)(1 + x )(1 + x ) (1 + x ) + .\]
1 x \cdots 1 x\par
- -\par
Deduce further that\par
\infty\par
2r\par
\[ln(1 x) = ln 1 + x ,\]
- -\par
\[r=0\]
X\par
and hence that\par
\[1 1 2x 4x3\]
\[= + + + .\]
\[1 x 1 + x 1 + x^2 1 + x^4 \cdots\]
-\par
\textbf{(ii)} Show that\par
\[1 + 2x 1 2x 2x 4x3 4x3 8x7\]
\[= - + - + - + .\]
\[1 + x + x^2 1 x + x^2 1 x^2 + x^4 1 x^4 + x^8 \cdots\]
- - -\par
\end{flushleft}
\subsection*{Page 2}
\begin{flushleft}
3 It is given that the two curves\par
2 2\par
\[y = 4 x and mx = k y ,\]
- -\par
\[where m > 0, touch exactly once.\]
\textbf{(i)} In each of the following four cases, sketch the two curves on a single diagram, noting\par
the coordinates of any intersections with the axes:\par
\[(a) k < 0 ;\]
\[(b) 0 < k < 16, k/m < 2 ;\]
\[(c) k > 16, k/m > 2 ;\]
\[(d) k > 16, k/m < 2 .\]
\textbf{(ii)} Now set m = 12.\par
Show that the x-coordinate of any point at which the two curves meet satisfies\par
4 2\par
\[x 8x + 12x + 16 k = 0 .\]
- -\par
Let a be the value of x at the point where the curves touch. Show that a satisfies\par
\[a 4a + 3 = 0\]
-\par
and hence find the three possible values of a.\par
Derive also the equation\par
\[k = 4a + 9a + 16 .\]
-\par
Which of the four sketches in part (i) arise?\par
\bigskip\textbf{Question 4.} (i) Show that\par
\infty\par
n + 1\par
\[= 2e 1\]
n! -\par
\[n=1\]
X\par
and\par
\infty\par
(n + 1)2\par
\[= 5e 1 .\]
n! -\par
\[n=1\]
X\par
\infty\par
(2n 1)3\par
Sum the series - .\par
n!\par
\[n=1\]
X\par
\[\infty (n^2 + 1)2 -n\]
\textbf{(ii)} Sum the series , giving your answer in terms of natural logarithms.\par
(n + 1)(n + 2)\par
\[n=0\]
X\par
\end{flushleft}
\subsection*{Page 3}
\begin{flushleft}
\bigskip\textbf{Question 5.} (i) The point with coordinates (a, b), where a and b are rational numbers, is called:\par
an integer rational point if both a and b are integers;\par
a non-integer rational point if neither a nor b is an integer.\par
(a) Write down an integer rational point and a non-integer rational point on the circle\par
\[x^2 + y^2 = 1.\]
\[(b) Write down an integer rational point on the circle x^2 + y^2 = 2. Simplify\]
\[(cos θ + \sqrt\{m\} sin θ) 2 + (sin θ \sqrt\{m\} cos θ) 2\]
-\par
\[and hence obtain a non-integer rational point on the circle x^2 + y^2 = 2 .\]
\textbf{(ii)} The point with coordinates (p+\sqrt2 q , r+\sqrt2 s), where p, q, r and s are rational numbers,\par
is called:\par
an integer 2-rational point if all of p, q, r and s are integers;\par
a non-integer 2-rational point if none of p, q, r and s is an integer.\par
(a) Write down an integer 2-rational point, and obtain a non-integer 2-rational point,\par
\[on the circle x^2 + y^2 = 3 .\]
\[(b) Obtain a non-integer 2-rational point on the circle x^2 + y^2 = 11 .\]
\[(c) Obtain a non-integer 2-rational point on the hyperbola x^2 y^2 = 7.\]
-\par
\bigskip\textbf{Question 6.} Let x + iy be a root of the quadratic equation z2 + pz + 1 = 0, where p is a real number. Show\par
\[that x^2 y^2 + px + 1 = 0 and (2x + p)y = 0. Show further that\]
-\par
\[either p = 2x or p = (x + 1)/x with x = 0 .\]
- - 6\par
Hence show that the set of points in the Argand diagram that can (as p varies) represent roots\par
of the quadratic equation consists of the real axis with one point missing and a circle. This\par
set of points is called the root locus of the quadratic equation.\par
Obtain and sketch in the Argand diagram the root locus of the equation\par
\[pz + z + 1 = 0\]
and the root locus of the equation\par
2 2\par
\[pz + p z + 2 = 0 .\]
\end{flushleft}
\subsection*{Page 4}
\begin{flushleft}
\bigskip\textbf{Question 7.} A pain-killing drug is injected into the bloodstream. It then diffuses into the brain, where it\par
is absorbed. The quantities at time t of the drug in the blood and the brain respectively are\par
y(t) and z(t). These satisfy\par
\[y˙ = 2(y z) , z˙ = y˙ 3z ,\]
- - - -\par
where the dot denotes differentiation with respect to t.\par
Obtain a second order differential equation for y and hence derive the solution\par
-t -6t 1 -t -6t\par
\[y = Ae + Be , z = Ae 2Be ,\]
2 -\par
where A and B are arbitrary constants.\par
\textbf{(i)} Obtain the solution that satisfies z(0) = 0 and y(0) = 5. The quantity of the drug in\par
\[the brain for this solution is denoted by z^1(t).\]
\textbf{(ii)} Obtain the solution that satisfies z(0) = z(1) = c, where c is a given constant. The\par
\[quantity of the drug in the brain for this solution is denoted by z^2(t).\]
\textbf{(iii)} Show that for 0 6 t 6 1,\par
\[z^2(t) = z^1(t n) ,\]
-\par
\[n=-\infty\]
X\par
provided c takes a particular value that you should find.\par
\bigskip\textbf{Question 8.} The sequence F0, F1, F2, . . . is defined by F0 = 0, F1 = 1 and, for n > 0,\par
\[Fn+2 = Fn+1 + Fn .\]
\textbf{(i)} Show that F0F3 F1F2 = F2F5 F3F4 .\par
- -\par
\textbf{(ii)} Find the values of FnFn+3 Fn+1Fn+2 in the two cases that arise.\par
-\par
\textbf{(iii)} Prove that, for r = 1, 2, 3, . . . ,\par
1 1 1\par
\[arctan = arctan + arctan\]
\[F2r F2r+1 F2r+2\]
\lambda \mu \lambda \mu \lambda \mu\par
and hence evaluate the following sum (which you may assume converges):\par
\infty\par
arctan .\par
\[F2r+1\]
\[r=1 \lambda \mu\]
X\par
\end{flushleft}
\subsection*{Page 5}
\begin{flushleft}
\bigskip\textbf{Section B: Mechanics}\par
\bigskip\textbf{Question 9.} A pulley consists of a disc of radius r with centre O and a light thin axle through O perpen-\par
dicular to the plane of the disc. The disc is non-uniform, its mass is M and its centre of mass\par
is at O. The axle is fixed and horizontal.\par
\[Two particles, of masses m^1 and m^2 where m^1 > m^2, are connected by a light inextensible\]
string which passes over the pulley. The contact between the string and the pulley is rough\par
enough to prevent the string sliding. The pulley turns and the vertical force on the axle is\par
found, by measurement, to be P + M g.\par
\textbf{(i)} The moment of inertia of the pulley about its axle is calculated assuming that the pulley\par
rotates without friction about its axle. Show that the calculated value is\par
\[((m^1 + m^2)P 4m1m2g)r^2\]
- . ( )\par
\[(m^1 + m^2)g P ∗\]
-\par
\textbf{(ii)} Instead, the moment of inertia of the pulley about its axle is calculated assuming that a\par
couple of magnitude C due to friction acts on the axle of the pulley. Determine whether\par
this calculated value is greater or smaller than ( ).\par
∗\par
\[Show that C < (m^1 m^2)rg.\]
-\par
\bigskip\textbf{Question 10.} A small ring of mass m is free to slide without friction on a hoop of radius a. The hoop is\par
fixed in a vertical plane. The ring is connected by a light elastic string of natural length a to\par
the highest point of the hoop. The ring is initially at rest at the lowest point of the hoop and\par
is then slightly displaced. In the subsequent motion the angle of the string to the downward\par
vertical is φ. Given that the ring first comes to rest just as the string becomes slack, find an\par
expression for the modulus of elasticity of the string in terms of m and g.\par
Show that, throughout the motion, the magnitude R of the reaction between the ring and the\par
hoop is given by\par
\[R = (12 cos φ 15 cos φ + 5)mg\]
-\par
and that R is non-zero throughout the motion.\par
\end{flushleft}
\subsection*{Page 6}
\begin{flushleft}
11 One end of a thin heavy uniform inextensible perfectly flexible rope of length 2L and mass\par
2M is attached to a fixed point P . A particle of mass m is attached to the other end. Initially,\par
the particle is held at P and the rope hangs vertically in a loop below P . The particle is then\par
released so that it and a section of the rope (of decreasing length) fall vertically as shown in\par
the diagram.\par
P\par
x\par
L + x\par
You may assume that each point on the moving section of the rope falls at the same speed as\par
the particle. Given that energy is conserved, show that, when the particle has fallen a distance\par
\[x (where x < 2L), its speed v is given by\]
2gx mL + M L M x)\par
\[v 2 = - 4 .\]
mL + M L M x\par
- 2\par
Hence show that the acceleration of the particle is\par
M gx mL + M L M x\par
g + - 4 .\par
1 2\par
2 mL + M L - 2 M x\par
Deduce that the acceleration of the particle after it is released is greater than g.\par
\end{flushleft}
\subsection*{Page 7}
\begin{flushleft}
\bigskip\textbf{Section C: Probability and Statistics}\par
\bigskip\textbf{Question 12.} (i) A point P lies in an equilateral triangle ABC of height 1. The perpendicular distances\par
\[from P to the sides AB, BC and CA are x^1, x^2 and x^3, respectively. By considering\]
\[the areas of triangles with one vertex at P , show that x^1 + x^2 + x^3 = 1.\]
Suppose now that P is placed at random in the equilateral triangle (so that the prob-\par
ability of it lying in any given region of the triangle is proportional to the area of that\par
region). The perpendicular distances from P to the sides AB, BC and CA are random\par
\[variables X^1, X^2 and X^3, respectively. In the case X^1 = min(X^1, X^2, X^3), give a sketch\]
showing the region of the triangle in which P lies.\par
\[Let X = min(X^1, X^2, X^3). Show that the probability density function for X is given by\]
6(1 3x) 0 6 x 6 1 ,\par
\[f(x) = - 3\]
(0 otherwise .\par
Find the expected value of X.\par
\textbf{(ii)} A point is chosen at random in a regular tetrahedron of height 1. Find the expected\par
value of the distance from the point to the closest face.\par
[The volume of a tetrahedron is area of base height and its centroid is a distance\par
\[1 3 \times \times\]
height from the base.]\par
\[4 \times\]
\bigskip\textbf{Question 13.} (i) The random variable Z has a Normal distribution with mean 0 and variance 1. Show\par
\[that the expectation of Z given that a < Z < b is\]
\[exp( 1 a^2) exp( 1 b^2)\]
- 2 - - 2 ,\par
\[\sqrt\{2\}\pi Φ(b) Φ(a)\]
-\par
where Φ denotes the cumulative distribution function for Z.\par
\textbf{(ii)} The random variable X has a Normal distribution with mean µ and variance σ2. Show\par
that\par
\[E(X X > 0) = µ + σE(Z Z > µ/σ).\]
| | -\par
Hence, or otherwise, show that the expectation, m, of X is given by\par
| |\par
1 2 2\par
\[m = µ 1 2Φ( µ/σ) + σ 2/\pi exp( µ /σ ) .\]
- - - 2\par
p\par
Obtain an expression for the variance of X in terms of µ, σ and m.\par
| |\par
\end{flushleft}
\end{document}
